\chapter{Notation, Open Problems, and Numerical Index}
\label{app:gaussian-backmatter}

\section{Notation}

This supplement uses the symbols of the main manuscript (\Cref{tab:notation}) with their meanings
there. \Cref{tab:gaussian-notation} lists the symbols that recur across the chapters of this
supplement; a symbol used inside one chapter is declared at its opening. The word pullback means
restriction of a tensor along an injective map of parameters, as in
\Cref{prop:ig-pullback-vs-pushforward}.

{\footnotesize
\renewcommand{\arraystretch}{1.12}
\begin{longtable}{@{}>{\raggedright\arraybackslash}p{0.29\textwidth}
>{\raggedright\arraybackslash}p{0.46\textwidth}
>{\raggedright\arraybackslash}p{0.19\textwidth}@{}}
\caption{Notation of Supplement G.}
\label{tab:gaussian-notation}\\
\toprule
\textbf{Symbol} & \textbf{Meaning} & \textbf{Introduced} \\
\midrule
\endfirsthead
\toprule
\textbf{Symbol} & \textbf{Meaning} & \textbf{Introduced} \\
\midrule
\endhead
\midrule
\multicolumn{3}{r@{}}{\emph{continued on the next page}} \\
\endfoot
\bottomrule
\endlastfoot
$K$, $d_m$; $n=NK$ & dimensions of the belief and model sample fibers, not assumed equal; dimension of the stacked belief latent & \cref{ch:gaussian-introduction,sec:gauss-interaction-family} \\
$\Gamma_b=\rho_b(G)$, $\Gamma_m=\rho_m(G)$ & represented frame groups, with the specialization $\Gamma_b=\GL^+(K)$ where it is named & \cref{hyp:gen-lg-frame-group} \\
$\nu$, $\mathsf T$, $\lambda$, $\mathsf A(\lambda)$, $\boldsymbol\mu_\lambda$ & reference measure, statistic, natural parameter, log partition, and mean parameter of a canonical family & \cref{def:exp-canonical-family} \\
$\mathcal W$ & closed convex operator cone preserved by the algebraic maps under study & \cref{def:exp-operator-layer} \\
$(\mathsf h,\Lambda)$; $\mu=\Lambda^{-1}\mathsf h$, $\Sigma=\Lambda^{-1}$; $\mathsf A(\mathsf h,\Lambda)$ & information vector and precision of a Gaussian law; its mean and covariance; its log normalizer & \cref{sec:gauss-information-form} \\
$F_i^b$, $F_i^m$, $\mathsf B_i$, $\mathsf C_i^b$, $\mathsf C_i^m$ & gains, model-to-belief bridge, and observation matrices of the directed linear-Gaussian forest & \cref{def:gen-lg-forest} \\
$\xi_i^b$, $\xi_i^m$, $\xi_i^o$; $\Sigma_i^b$, $\Sigma_i^m$, $\Sigma_i^o$; $\operatorname{par}(i)$ & offsets and receiving covariances of the belief, model, and observation factors; the parent of a nonroot agent & \cref{def:gen-lg-forest} \\
$\Omega_i$ & transition $F_i\Theta_{i\operatorname{par}(i)}$ of a nonroot agent & \cref{prop:gauss-assembled-precision-blocks} \\
$\rho^b_i=\rho_b(g_i)$, $\rho^m_i=\rho_m(g_i)$ & represented passive re-expressions of agent $i$, the maps $R_i^b$ and $R_i^m$ of the main manuscript & \cref{prop:gen-lg-kernel-covariance} \\
$\mathcal I(V)$; $A_i$, $W_{ij}=W_{ji}$, $L$ & interaction family on the agents; self term, edge weight, and matrix-weighted Laplacian part, with $\Lambda_{ii}=A_i+\sum_{j\neq i}W_{ij}$ and $\Lambda_{ij}=-W_{ij}$ & \cref{def:gauss-interaction-family} \\
$O_i$ & observation precision added to the self term of agent $i$ & \cref{prop:gauss-gain-absorption} \\
$z_i=\rho_b(U_i)^{-1}\mu_i$ & mean of agent $i$ in link-trivializing coordinates & \cref{prop:gauss-vector-residuals} \\
$\mathbf P=\operatorname{diag}(\mathbf P_1,\ldots,\mathbf P_N)$, $\mathbf P_i\in\Gamma_b$ & block reframing from new to old coordinates, under which a precision becomes $\mathbf P^\top\Lambda\mathbf P$ & \cref{def:ig-block-reframing} \\
$S$, $n_c$; $n_I$ & aggregation matrix of a hard partition into $n_c$ clusters; size of the cluster $I$ & \cref{prop:cg-gaussian-aggregation} \\
$P_\Lambda=(S^\top\Lambda S)^{-1}S^\top\Lambda$ & coarse coordinate map carrying the Galerkin law; $SP_\Lambda$ is the $\Lambda$-orthogonal projector onto cluster-constant configurations & \cref{prop:ig-galerkin-pushforward} \\
$\widehat\Lambda$, $\widehat{\mathsf h}$ & precision and information vector of the joint law of $(o,Y)$ given the structural data & \cref{prop:restrict-conditional-precision-observation-free} \\
$\mathsf E$ & energy of a law proportional to $e^{-\mathsf E}$ & \cref{def:cg-graph-exponential-energy} \\
$\mathcal W_H$, $M_0$ & congruence-diagonal cone with frame $H$, $HH^\top=M_0$; an SPD order unit of a cone & \cref{def:cg-congruence-diagonal-cone} \\
$\boldsymbol\Xi_\ell$, $\mathcal A_S$, $\zeta_\ell$ & operator parameters $((A_i^{(\ell)}),(W_{ij}^{(\ell)}))$ at depth $\ell$, the unrescaled aggregation map, and the block rescaling & \cref{def:grg-operator-arrow} \\
$L^2_0$, $b$; $\mathscr I_b$, $\mathscr L_b$ & centered square-integrable scores of a law, block size; score replication lift and score operator of a block of $b$ copies & \cref{def:rg-dqm-score,thm:rg-score-pushforward-defect} \\
$\mathrm{He}_k$, $\widehat{\mathrm{He}}_k=\mathrm{He}_k/\sqrt{k!}$ & probabilists' Hermite polynomial and the normalized Hermite basis & \cref{thm:rg-gaussian-hermite-spectrum} \\
\end{longtable}
}

\section{Open problems}
\label{sec:grg-obligations}

\Cref{tab:gaussian-open} lists the open problems of this supplement, each with the statement at which
it arises and a condition that would settle it. The first three arise across chapters and are stated
here. The Gaussian chapters also bear on open problems of the main manuscript, among them the choice
of a nondegenerate recognition structure (\Cref{open:restrict-intrinsic-criterion}) and of a natural
partition (\Cref{open:cg-natural-selector}). The Gaussian coarse families also meet the established
obstruction to a continuous extension of balanced cut links (\Cref{open:cg-cut-compression}).

\openproblemheading{Bayesian renormalization and Gaussian aggregation}{open:grg-bayesian-bridge}
Aggregation of interaction parameters is a positive semigroup that sums cut weights and self terms
(\Cref{prop:grg-aggregation-semigroup}), and a rescaled step on the interior maps normalized Gaussian
laws to normalized Gaussian laws (\Cref{prop:grg-interior-law-step}). In its late-data regime Bayesian
renormalization works with a Gaussian posterior whose covariance is the scale parameter times the
inverse Fisher information (\Cref{sec:rg-bayesian-comparison}). Neither this supplement nor the main
manuscript relates the two constructions, and in particular neither shows that they run in opposite
directions; similar Gaussian formulas do not supply such a relation. The problem is settled by a common state space carrying both
constructions, an explicit map between the Bayesian scale parameter and a declared block ratio with
its rescaling, the Fisher tensors transported along that map, and a proof that a transported
functional is monotone along both constructions, or by a Gaussian hierarchy in a stated class for
which no such functional exists.

\openproblemheading{Transfer of a model-space Fisher scale}{open:ig-model-space-fisher-transfer}
At a fixed precision \(\Lambda\) the Fisher metric of the Gaussian mean family is \(\Lambda\)
(\Cref{prop:ig-fisher-moment-chart}), and restriction and marginalization give the two metrics compared
in \Cref{prop:ig-pullback-vs-pushforward}. The information scale of Bayesian renormalization is built
from a Fisher metric on the parameters of a generative model and is written with its inverse
(\Cref{sec:rg-bayesian-comparison}). Whether that construction applies to \(\Lambda\) read as the
Fisher metric of a recognition mean is not settled by the common name: a coarse recognition mean is not
a parameter of the generative law, and the construction does not treat a singular metric. The problem is settled by a declaration that
makes the coarse recognition mean a parameter of the generative law fitted against the evidence bound,
a proof that the Fisher metric in that parameter is nondegenerate on the declared family, and then the
transferred construction with its proof, or a declared family in which that metric is singular.

\openproblemheading{Stochastic inversion of aggregation}{open:grg-stochastic-inverse}
Hard aggregation annihilates every edge internal to a cluster
(\Cref{prop:cg-gaussian-aggregation}), so the aggregation map \(\mathcal A_S\) on operator parameters
is not injective (\Cref{prop:grg-aggregation-semigroup}). A refinement kernel that assigns to each
coarse parameter tuple a law on its preimage under \(\mathcal A_S\) assigns the same law to every fine
tuple with that image, so it redraws fine parameters compatible with the coarse ones, and on a
preimage with more than one point it returns the realized fine tuple almost surely for at most one of
them. What is open is a stochastic inverse of a hierarchy of aggregations.
It is settled by priors for the internal edge weights discarded at every step, refinement kernels that
are right inverses of the aggregation maps and compose across depths, the chart data that identify the
levels, and a projective-limit theorem for the refined laws, or by a proof that no compatible family
exists for a stated class of hierarchies.

{\footnotesize
\renewcommand{\arraystretch}{1.12}
\begin{longtable}{@{}>{\raggedright\arraybackslash}p{0.25\textwidth}
>{\raggedright\arraybackslash}p{0.17\textwidth}
>{\raggedright\arraybackslash}p{0.52\textwidth}@{}}
\caption{Open problems of Supplement G.}
\label{tab:gaussian-open}\\
\toprule
\textbf{Item} & \textbf{Locator} & \textbf{Settling condition} \\
\midrule
\endfirsthead
\toprule
\textbf{Item} & \textbf{Locator} & \textbf{Settling condition} \\
\midrule
\endhead
\midrule
\multicolumn{3}{r@{}}{\emph{continued on the next page}} \\
\endfoot
\bottomrule
\endlastfoot
Bayesian renormalization and Gaussian aggregation & \cref{open:grg-bayesian-bridge} & A common state space, a map between the Bayesian scale parameter and a declared block ratio with its rescaling, transported Fisher tensors, and a transported functional proved monotone along both constructions; or a Gaussian hierarchy in a stated class with no such functional. \\
Transfer of a model-space Fisher scale & \cref{open:ig-model-space-fisher-transfer} & A declaration making the coarse recognition mean a parameter of the generative law fitted against the evidence bound, nondegeneracy of the Fisher metric in that parameter, and the transferred construction with proof; or a declared family in which that metric is singular. \\
Stochastic inversion of aggregation & \cref{open:grg-stochastic-inverse} & Priors for the discarded internal edge weights, refinement kernels that are right inverses of the aggregation maps and compose across depths, chart data, and a projective-limit theorem; or a proof that no compatible family exists for a stated class of hierarchies. \\
Scaling of a retained mass sector & \cref{open:exp-mass-sector} & For a given scale map, an explicit induced law on the mass coordinate together with a proof, under stated hypotheses on the map, that its orbit stays in the natural domain or reaches its boundary. \\
Commutation of the information projection & \cref{open:exp-projection-commutation} & An intertwining theorem with hypotheses on a gauge action or a coarse parameter map under which the affine information projection commutes with it, or a counterexample within a regular minimal family. \\
Monotonicity along aggregation with varying precision & \cref{open:grg-galerkin-monotonicity} & A Lyapunov functional with a proof of its decrease for families on which $\operatorname{range}(\Lambda S)$ varies; a declared quotient, pin, or mass repair of the boundary member of the bi-additive family, with its own reference measure, normalizer, and pushforward identity; and a projective-limit theorem for composed kernels at unbounded depth. \\
Kron classification beyond finite encodings & \cref{open:cg-kron-classification-beyond-finite} & For a stated family outside finite semialgebraic or quotient encodings, a declaration of its domain and compatibility conditions with a proof of the requested invariance, limit, or global construction; or a reduction of the family to a finite encoding covered by \Cref{thm:cg-semialgebraic-kron-classification,thm:cg-variable-support-classification}. \\
Robustness of the star schedule & \cref{open:obs-star-update-robustness} & A convergence theorem, with a rate, for a declared class of delayed, noisy, asynchronous, or inexact updates of the anchored star, or a counterexample within such a class. \\
\end{longtable}
}

\section{Numerical index}

\Cref{tab:gaussian-numerical} indexes the computations recorded for this supplement, in the columns
fixed for every numerical index. No statement of this supplement has the status NUMERICAL. Each row
is a regression check of established statements, which rest on their proofs; the checks corroborate
the declared finite instances and prove nothing beyond them.

{\footnotesize
\renewcommand{\arraystretch}{1.12}
\begin{longtable}{@{}>{\raggedright\arraybackslash}p{0.16\textwidth}
>{\raggedright\arraybackslash}p{0.17\textwidth}
>{\raggedright\arraybackslash}p{0.28\textwidth}
>{\raggedright\arraybackslash}p{0.17\textwidth}
>{\raggedright\arraybackslash}p{0.11\textwidth}@{}}
\caption{Numerical index of Supplement G.}
\label{tab:gaussian-numerical}\\
\toprule
\textbf{Claim} & \textbf{Check} & \textbf{What is tested} & \textbf{Control} & \textbf{Seed} \\
\midrule
\endfirsthead
\toprule
\textbf{Claim} & \textbf{Check} & \textbf{What is tested} & \textbf{Control} & \textbf{Seed} \\
\midrule
\endhead
\bottomrule
\endlastfoot
\cref{prop:cg-mean-tie-cost}, \eqref{eq:cg-factorization-gap} & \texttt{CHK-CG-FACTOR-\allowbreak GAP-\allowbreak STRESS-\allowbreak 3138} & The factorization gap computed in binary64 through canonical correlations, against an evaluation of the same stored inputs at 100 decimal digits, on 3,138 deterministic cases of dimensions 2 to 16 with nominal condition numbers from 1 to $10^{14}$. & 18 designated inputs evaluated exactly at 100 decimal digits. & 20260803 \\
\cref{prop:obs-anchor-coercivity,cor:obs-flat-fold-singular}, \eqref{eq:obs-precision-addition} & \texttt{CHK-OBS-\allowbreak STAR-\allowbreak FOLD-\allowbreak NEW-\allowbreak PROTOCOL} & Anchored-star and flat reciprocal-pair instances: the apex precision-addition identity and the contrast of the two forms on transport-consistent configurations. & Flat unanchored pairs as comparison instances. & 20260728 \\
\cref{prop:obs-holonomy-determinant-factorization,cor:obs-holonomy-kernel-shrinkage,prop:obs-normalizer-link-dependence} & \texttt{CHK-OBS-\allowbreak HOLONOMY-\allowbreak DET-\allowbreak KERNEL}, \texttt{CHK-OBS-\allowbreak NORMALIZER-\allowbreak WITNESS} & Reciprocal pairs with explicit transports and covariances: the determinant identity, the prescribed dimensions of the kernel, and the normalizer witness. & A normalized anchored instance. & 20260728 \\
\end{longtable}
}
